\documentclass[UTF8]{ctexart}

\usepackage{amscd,amsthm,amsfonts,amssymb,amsmath}
\usepackage[colorlinks=true]{hyperref}
\usepackage[all]{xy}
\usepackage{mathrsfs}
\usepackage{tikz,mathpazo}
\usepackage{bm}
\usepackage{geometry}
\usepackage{xcolor}
\usepackage{eso-pic}
%\usepackage{CJK}
\usepackage{arydshln}
\usepackage{mathptmx}
\usepackage[11pt]{moresize}

\newcommand\bs[1]{\boldsymbol{ #1}}

\newenvironment{bmatriX}[1]{\left[\hskip -\arraycolsep \begin{array}{#1}}{\end{array}\hskip -\arraycolsep\right]}
	
	\newtheorem{thm}{定理}
	\newtheorem{cor}[thm]{推论}
	\newtheorem{lem}[thm]{引理}
	\newtheorem{prop}[thm]{命题}
	\newtheorem{defn}[thm]{定义}
	\newtheorem{ques}[thm]{问题}
	\newtheorem{eg}[thm]{例题}
	\newtheorem{ex}[thm]{习题}
	\newtheorem{rmk}[thm]{注释}
		\newtheorem{ntn}[thm]{记号}

\newif\ifhideproofs
\hideproofstrue
\newif\iffinalver
%\finalvertrue        
\ifhideproofs
\usepackage[final]{changes}
\definechangesauthor[name=Error, color=red]{err}
\usepackage{environ}
\NewEnviron{hide}{}
\let\proof\hide
\let\endproof\endhide
\else
\iffinalver
\usepackage[final]{changes}
\else
\usepackage{changes}
\fi
\definechangesauthor[name=Error, color=red]{err}
\fi

		
\begin{document}


\title{习题课材料（十一）}

\date{}

\maketitle

\vspace{-1cm}



\noindent{\textbf{我们用$\mathbb{F}$表示一个给定的数域。}}

\begin{ex}
给定$A = \begin{bmatrix}
1 & 2 & 1 \\
1 & 1 & 0 \\
0 & 1 & 1
\end{bmatrix}$和$\mathbb{R}^{3 \times 3}$上的线性变换$f: X \to AX$。分别求$\mathcal{N}(f)$和$\mathcal{R}(f)$的维数和一组基。
\end{ex}
\begin{proof}[解]
%由定义$\mathcal{N}(f) = \{X \in \mathbb{R}^{3 \times 3} \mid AX = 0\}$
定义线性空间的同构$\Phi: \mathbb{R}^{3 \times 3} \to \mathbb{R}^{3} \oplus \mathbb{R}^{3} \oplus \mathbb{R}^{3}: [\alpha_1, \alpha_2, \alpha_3] \to (\alpha_1, \alpha_2, \alpha_3)$，即从列向量的角度分解矩阵。注意到如下的交换图：
\begin{equation}
\xymatrix{
[\alpha_1, \alpha_2, \alpha_3] \in \mathbb{R}^{3 \times 3}\ar[r]^{\Phi}\ar[d]^{f}& (\alpha_1, \alpha_2, \alpha_3) \in \mathbb{R}^{3} \oplus \mathbb{R}^{3} \oplus \mathbb{R}^{3}\ar[d]\\
A[\alpha_1, \alpha_2, \alpha_3] \in \mathbb{R}^{3 \times 3}\ar[r]^{\Phi}& (A\alpha_1, A\alpha_2, A\alpha_3) \in \mathbb{R}^{3} \oplus \mathbb{R}^{3} \oplus \mathbb{R}^{3}.\\
}
\end{equation}
右侧垂直映射记为$(f_1, f_2, f_3)$，每个$f_i$都是左乘矩阵$A$。验证:
\

$\mathcal{N}(f) = \Phi^{-1}(\mathcal{N}((f_1, f_2, f_3)))$以及$\mathcal{R}(f) = \Phi^{-1}(\mathcal{R}((f_1, f_2, f_3)))$。
\

注意到$\mathcal{N}(f_1) = \mathcal{N}(f_2) = \mathcal{N}(f_3) = \mathcal{N}(A) = \mathbb{R}\begin{bmatrix} 
1 \\
-1 \\
1
\end{bmatrix} = \{\begin{bmatrix} 
a \\
-a \\
a
\end{bmatrix} \mid a \in \mathbb{R}\}$以及
$\mathcal{R}(f_1) = \mathcal{R}(f_2) = \mathcal{R}(f_3) = \mathcal{R}(A) = \mathrm{Span} \langle \begin{bmatrix} 
1 \\
1 \\
0
\end{bmatrix}, \begin{bmatrix} 
1 \\
0 \\
1
\end{bmatrix}\rangle = \{\begin{bmatrix} 
b+c \\
b \\
c
\end{bmatrix} \mid b,c \in \mathbb{R} \}$，因此
\[
\mathcal{N}(f) = \{
\begin{bmatrix} 
a_1 & a_2 & a_3\\
-a_1 & -a_2 & -a_3\\
a_1 & a_2 & a_3
\end{bmatrix} \mid a_1, a_2, a_3 \in \mathbb{R} \}= \mathrm{Span} \langle E_{11} - E_{21}+E_{31}, E_{12} - E_{22}+E_{32}, E_{13} - E_{23}+E_{33}\rangle
\]
以及
\begin{equation} \nonumber
\begin{split}
&\mathcal{R}(f) = \{
\begin{bmatrix} 
b_1+c_1 & b_2+c_2 & b_3+c_3\\
b_1 & b_2 & b_3\\
c_1 & c_2 & c_3
\end{bmatrix} \mid b_1, b_2, b_3, c_1, c_2, c_3 \in \mathbb{R}
\} \\
= &\mathrm{Span} \langle E_{11}+E_{21}, E_{11}+E_{31}, E_{12}+E_{22}, E_{12}+E_{32},E_{13}+E_{23}, E_{13}+E_{33}\rangle.
\end{split}
\end{equation}
于是$\dim \mathcal{N}(f) = 3, \dim \mathcal{R}(f) = 6$。
\end{proof}
\begin{rmk}
此题也可以直接计算$f$在基$E_{ij}, 1 \leq i,j \leq 3$下的表示矩阵。
\end{rmk}




\begin{ex}
设$\mathbb{Q}[i] = \{a+bi \mid a, b \in \mathbb{Q} \}$，其中$i^2 = -1$。
\begin{enumerate}
\item
证明$\mathbb{Q}[i]$关于数的加法和数乘构成 $\mathbb{Q}$上的线性空间。
\item
计算$\mathbb{Q}[i]$作为$\mathbb{Q}$上线性空间的维数。
\end{enumerate}
\end{ex}
\begin{proof}[解]
\begin{enumerate}
\item
验证线性空间的八条性质，细节略。
\item
先证明$\{1, i\}$构成$\mathbb{Q}[i]$的一组基。首先，由定义，$\mathbb{Q}[i]$中的任意元素可以写成$1$和$i$的线性组合。其次，若存在$m, n \in \mathbb{Q}$满足
$m + ni = 0,$
则
\[
0 = (m+ni)(m-ni) = m^2 + n^2,
\]
因此$m = n = 0$，换言之$1$和$i$线性无关。所以$\dim_{\mathbb{Q}}\mathbb{Q}[i] = 2$。
\end{enumerate}
\end{proof}



\begin{ex}
设$V$是数域$\mathbb{F}$上的$n$维线性空间，$V_1$和$V_2$是$V$的两个子空间，满足$\dim V_1 + \dim V_2 = \dim V$。请问是否存在$V$上的线性变换$\sigma$满足$\mathcal{N}(\sigma) = V_1$且 $\mathcal{R}(\sigma) = V_2$？
\end{ex}
\begin{proof}[解]
存在。设向量组$\{\alpha_1, \cdots, \alpha_r\}$与$\{\beta_{r+1}, \cdots, \beta_{n}\}$分别构成$V_1$与$V_2$的一组基。将$\{\alpha_1, \cdots, \alpha_r\}$扩充为$V$的一组基$\{\alpha_1, \cdots, \alpha_r, \alpha_{r+1}, \cdots, \alpha_n\}$。定义$V$上的线性变换$\sigma$：
\[
\sigma(\alpha_i) = \begin{cases}
0, & 1 \leq i \leq r; \\
\beta_i, & r+1 \leq i \leq n.
\end{cases}
\]
验证$\sigma$满足题目中的条件。
\end{proof}



\begin{ex}
定义$\mathbb{F}[x]$上的变换$\mathbf{A}(f(x)) = xf(x), \forall f(x) \in \mathbb{F}[x]$。
\begin{enumerate}
\item
验证$\mathbf{A}$是$\mathbb{F}[x]$上的一个线性变换。
\item
设$\mathbf{D}$是求导算子，证明$\mathbf{D}\mathbf{A} - \mathbf{A}\mathbf{D} = \mathbf{I}$。
\end{enumerate}
\end{ex}
\begin{proof}
\begin{enumerate}
\item
由线性变换的定义，需验证，$\forall, \alpha, \beta \in \mathbb{F}, f(x), g(x) \in \mathbb{F}[x]$，
\[
\mathbf{A}(\alpha f(x)+\beta g(x)) = \alpha \mathbf{A}(f(x))  + \beta \mathbf{A}(g(x)). 
\]
由定义，上式左边为$\mathbf{A}(\alpha f(x)+\beta g(x)) = x(\alpha f(x)+\beta g(x)) = \alpha xf(x) + \beta xg(x) =$ 等式右边。所以$\mathbf{A}$是$\mathbb{F}[x]$上的一个线性变换。
\item
直接计算即可。
\begin{equation} \nonumber
\begin{split}
&(\mathbf{D}\mathbf{A} - \mathbf{A}\mathbf{D} )(f(x)) = \mathbf{D}\mathbf{A} f(x)- \mathbf{A}\mathbf{D} f(x) 
=  \mathbf{D}(xf(x)) - \mathbf{A}(f^{'}(x)) \\
=& (xf(x))^{'} - xf^{'}(x) = f(x) + xf^{'}(x) - xf^{'}(x) = f(x) = \mathbf{I} f(x). 
\end{split}
\end{equation}
\end{enumerate}
\end{proof}

\begin{ex} %%% 过渡矩阵
设 $a_{1} = \begin{bmatrix}
1 & 0 \\
0 & 0
\end{bmatrix}, a_{2} =\begin{bmatrix}
1 & 1 \\
0 & 0
\end{bmatrix}, a_{3} =\begin{bmatrix}
0 & 1 \\
1 & 0
\end{bmatrix}, a_{4} = \begin{bmatrix}
0 & 0 \\
1 & 1
\end{bmatrix}$;
\

$b_{1} = \begin{bmatrix}
1 & 0 \\
1 & 0
\end{bmatrix}, b_{2} =\begin{bmatrix}
0 & 1 \\
0 & 0
\end{bmatrix}, b_{3} =\begin{bmatrix}
0 & 0 \\
1 & 1
\end{bmatrix}, b_{4} = \begin{bmatrix}
0 & 0 \\
0 & 1
\end{bmatrix}$.
\begin{enumerate}
\item
验证$\{a_{i}, 1 \leq i \leq 4\}$和$\{b_{i}, 1 \leq i \leq 4\}$都组成$\mathbb{R}^{2 \times 2}$的一组基。
\item
计算$\{a_{i}, 1 \leq i \leq 4\}$到$\{b_{i}, 1 \leq i \leq 4\}$的过渡矩阵。
\end{enumerate}
\end{ex}
\begin{proof}[解]
\begin{enumerate}
\item
计算可知：
\[
[a_1, a_2, a_3, a_4] = [E_{11}, E_{12}, E_{21}, E_{22}] A = [E_{11}, E_{12}, E_{21}, E_{22}] \begin{bmatrix}
1 & 1 & 0 & 0  \\
0 & 1 & 1 & 0  \\
0 & 0 & 1 & 1 \\
0 & 0 & 0 & 1 
\end{bmatrix},
\]
\[
[b_1, b_2, b_3, b_4] = [E_{11}, E_{12}, E_{21}, E_{22}] B = [E_{11}, E_{12}, E_{21}, E_{22}] \begin{bmatrix}
1 & 0 & 0 & 0  \\
1 & 1 & 0 & 0  \\
0 & 0 & 1 & 0 \\
0 & 0 & 1 & 1 
\end{bmatrix}.
\]
由于\replaced{$A, B$}{$B, C$}可逆，故$\{a_{i}, 1 \leq i \leq 4\}$和$\{b_{i}, 1 \leq i \leq 4\}$都组成$\mathbb{R}^{2 \times 2}$的一组基。
\item
记$\{a_{i}, 1 \leq i \leq 4\}$到$\{b_{i}, 1 \leq i \leq 4\}$的过渡矩阵为$C$，即
\[
[b_1, b_2, b_3, b_4] = [a_1, a_2, a_3, a_4]C.
\]
于是：
\[
[E_{11}, E_{12}, E_{21}, E_{22}] B = [E_{11}, E_{12}, E_{21}, E_{22}] A C. 
\]
换句话说$C = A^{-1}B = \begin{bmatrix}
1 & 1 & 0 & 0  \\
0 & 1 & 1 & 0  \\
0 & 0 & 1 & 1 \\
0 & 0 & 0 & 1 
\end{bmatrix}^{-1} \begin{bmatrix}
1 & 0 & 0 & 0  \\
1 & 1 & 0 & 0  \\
0 & 0 & 1 & 0 \\
0 & 0 & 1 & 1 
\end{bmatrix} = \begin{bmatrix}
0 & -1 & 0 & -1  \\
1 & 1 & 0 & 1  \\
0 & 0 & 0 & -1 \\
0 & 0 & 1 & 1 
\end{bmatrix}$。
\end{enumerate}
\end{proof}
%%%%%%     1  -1  1  -1
%%%%%%         1   -1  1
%%%%%%             1  -1
%%%%%%                 1
\begin{ex} 
设$W$是数域$\mathbb{F}$上的$n$维线性空间，$V_1, V_2$和$V$都是$W$的子空间。假设$V_1 \subset V_2$，$V\cap V_1 = V\cap V_2$, $V+ V_1=V+V_2$，证明：
$V_1 = V_2$。
\end{ex}
\begin{proof}
由维数公式可知：
\[
\dim(V+ V_1) = \dim V + \dim V_1 - \dim(V \cap V_1), \dim(V+ V_2) = \dim V + \dim V_2 - \dim(V\cap V_2).
\]
由题设又知$\dim(V+V_1) = \dim(V+V_2)$，$\dim(V\cap V_1) = \dim(V\cap V_2)$，因此$\dim V_1 = \dim V_2$。结合$V_1 \subset V_2$，可知结论成立。
\end{proof}
\begin{ques}
题目中条件$V_1 \subset V_2$是否需要？
\end{ques}
%需要 V={x=y}, V_1={x=0}, V_2={y=0}


\begin{ex}
设$\mathbb{F}^{2 \times 2}$上的线性变换$\mathbf{R}_A: X \to XA$，其中$A = \begin{bmatrix} 
a & b \\
c & d
\end{bmatrix}$。求$\mathbf{R}_A$在基$E_{11}, E_{12}, E_{21}, E_{22}$下的矩阵。
\end{ex}
\begin{proof}[解]
\begin{equation} \nonumber
\mathbf{R}_A(E_{11}) = E_{11}A = \begin{bmatrix} 
1 & 0 \\
0 & 0
\end{bmatrix}\begin{bmatrix} 
a & b \\
c & d
\end{bmatrix} = \begin{bmatrix} 
a & b \\
0 & 0
\end{bmatrix} = aE_{11} + bE_{12}.
\end{equation}
同样计算可知：
\begin{equation} \nonumber
\mathbf{R}_A(E_{12}) = cE_{11} + dE_{12};
\end{equation}
\begin{equation} \nonumber
\mathbf{R}_A(E_{21}) = aE_{21} + bE_{22};
\end{equation}
\begin{equation} \nonumber
\mathbf{R}_A(E_{22}) = cE_{21} + dE_{22}.
\end{equation}
所以，我们有
\begin{equation} \nonumber
\mathbf{R}_A(E_{11}, E_{12}, E_{21}, E_{22}) = (E_{11}, E_{12}, E_{21}, E_{22})
\begin{bmatrix} 
a & c & 0 & 0 \\
b & d & 0 & 0 \\
0 & 0 & a & c \\
0 & 0 & b & d
\end{bmatrix}.
\end{equation}
\end{proof}




\begin{ex}
设$\mathcal{V}$是数域$\mathbb{F}$上的$n$维线性空间，将$\mathbb{F}$看作自身上的线性空间，而$\mathcal{V}$到$\mathbb{F}$的线性映射称为$\mathcal{V}$上的线性函数。令$\mathcal{V}^* = \mathrm{Hom}(\mathcal{V}, \mathbb{F})$，称为$\mathcal{V}$的对偶空间，证明$\mathcal{V}^{*}$和$\mathcal{V}$同构。
\end{ex}
\begin{proof}
固定$\mathcal{V}$的一组基，记为$\{e_1, \cdots, e_n\}$。定义映射$\varphi: \mathcal{V}^{*} \to \mathcal{V}$：
\[
f \to (e_1, e_2, \cdots, e_n)\begin{bmatrix}
f(e_1) \\
f(e_2) \\
\vdots \\
f(e_n)
\end{bmatrix}
\]
直接验证$\varphi$是线性同构即可。
\end{proof}
\begin{rmk}
作为推论，我们发现$\dim \mathcal{V}^{*} = \dim \mathcal{V}$。$\{e_1, \cdots, e_n\}$在$\varphi$下的原象构成了$\mathcal{V}^{*}$的一组基，称为$\mathcal{V}^{*}$关于$\{e_1, \cdots, e_n\}$的对偶基，一般记为$\{e^1, \cdots, e^n\}$。
\end{rmk}


\begin{ex}[$\heartsuit $]
给定$m$阶方阵$A_1$，$n$阶方阵$A_2$和$m \times n$阶矩阵$B$。证明：如果$A_1$和$A_2$没有相同的特征值，则
\begin{enumerate}
\item
关于$m \times n$矩阵$X$的方程，$A_1X = XA_2$只有平凡解。
\item
关于$m \times n$矩阵$X$的方程，$A_1X - XA_2 = B$只有唯一解。
\end{enumerate}
\end{ex}
\begin{proof}
\begin{enumerate}
\item
注意到：对于任意可逆矩阵$P$，$(P^{-1}A_1P)(P^{-1}XP) = (P^{-1}XP)(P^{-1}A_2P)$
只有平凡解当且仅当$A_1X = XA_2$只有平凡解。另一方面存在可逆矩阵$P$使得$P^{-1}A_2P$为上三角矩阵，所以不妨设$A_2$为上三角矩阵，其主对角元为$\lambda_1, \cdots, \lambda_n$。令$X = (x_1, \cdots, x_n)$为$A_1X = XA_2$的解。于是
\[
(A_1x_1, \cdots, A_1x_n) = A_1(x_1, \cdots, x_n) = (x_1, \cdots, x_n) A_2 =  (x_1, \cdots, x_n) \begin{bmatrix}
\lambda_1 & * & \cdots & * \\
0 & \lambda_2 & \cdots & * \\
0 & 0 & \cdots & \cdots \\
0 & 0 &  \cdots & \lambda_n
\end{bmatrix}.
\]
特别地，我们发现$A_1x_1 = \lambda_1 x_1$。因为$A_1$和$A_2$没有相同的特征值，故$A_1 - \lambda_1 \mathrm{I}_n$为可逆矩阵，所以$x_1 = 0$。再观察第二列，我们发现$A_1x_2 = \lambda_2 x_2 + * x_1 = \lambda_2 x_2$，同样利用$A_1$和$A_2$没有相同的特征值，$A_1 - \lambda_2 \mathrm{I}_n$为可逆矩阵，因此$x_2 = 0$。同理，我们可以依次证明$x_3 = \cdots = x_n = 0$，即$A_1X = XA_2$只有平凡解。
\added{
\begin{proof}
	可以证明对任意多项式$f(x)$,有$f(A)X=Xf(B)$，取$f$为$A$的特征多项式，则由$Caley-Hamilton$定理有$f(A)=0$，又由$f$的根为$A$的特征值，且$A, B$没有公共特征值，有$f(B)$可逆，于是可得$X=0$。
\end{proof}
}
\item
假设$m \times n$矩阵$X_1, X_2$都满足$A_1X - XA_2 = B$，即
\[
A_1X_1 - X_1A_2 = B, A_1X_2 - X_2A_2 = B. 
\]
两式相减得$A_1(X_1-X_2) - (X_1-X_2)A_2 = 0$。由第一问知，$X_1-X_2 = 0$，因此解唯一。

\added{
	由与第1问类似的讨论可知解的存在性，或将$X\mapsto A_1X-XA_2$看作$n$维矩阵线性空间上的线性变换，由第1问知是单射，故由维数公式知是满射。
}
\end{enumerate}
\end{proof}
\clearpage
\listofchanges
\end{document}
